\documentclass[conference]{IEEEtran}

\usepackage{cite}
\usepackage{amsmath,amssymb,amsfonts}
\usepackage{algorithmic}
\usepackage{graphicx}
\usepackage{textcomp}
\usepackage{xcolor}
\usepackage{booktabs}
\usepackage{array}
\usepackage{fvextra}
\usepackage{enumitem}
\usepackage{microtype}
\usepackage{hyperref}
\usepackage{tagging}
\usepackage{sc26repro}

\hypersetup{colorlinks=true,linkcolor=black,urlcolor=blue}
\setlist{nosep,leftmargin=*}

% SC26 final submission: explanation/example tags intentionally disabled.
% \usetag{explanation}
% \usetag{example}

\begin{document}

\twocolumn[%
{\begin{center}
\Huge
Appendix: Artifact Description/Artifact Evaluation
\end{center}}
]

\appendixAD

Paper ID: \texttt{pap142}


\section{Part 1: Overview of Contributions and Artifacts}

\subsection{Paper's Main Contributions}

\begin{itemize}
\item C1 - Target-centric execution. TempGNN introduces Target Dependency
Packets (TDPs) to express only the temporally valid states needed by a target
update, exposing branch-level parallelism while preserving ordered target
state commitment.
\item C2 - Dependency-driven construction. TempGNN combines temporal sampling,
dependency expansion, and data loading in an on-the-fly DDTC pipeline to
reduce target-irrelevant work and off-chip traffic.
\item C3 - Overlap-aware execution. TempGNN detects shared dependency packets
among concurrent targets and uses OATS/PHLE reuse while retaining
order-preserving target updates.
\item C4 - U280 evaluation. The paper evaluates the execution behavior,
performance, energy, ablations, and sensitivity of TempGNN across four TGNN
models and six temporal datasets, including comparisons with software and
FPGA baselines.
\end{itemize}

\subsection{Computational Artifacts}

Primary code repository:


\begin{Verbatim}[fontsize=\small,breaklines=true,breakanywhere=true]
https://github.com/TempGNN-Program/TempGNN
\end{Verbatim}

Reviewer-downloadable package:


\begin{Verbatim}[fontsize=\small,breaklines=true,breakanywhere=true]
https://github.com/TempGNN-Program/TempGNN/raw/main/ae_export/pap142_tempgnn_sc26_ae_u280.tgz
\end{Verbatim}

The version-specific DOI will be inserted before the SC26 artifact freeze on

August 25, 2026. GitHub alone is not the persistent archive; the frozen release

will also be deposited in Zenodo, Figshare, or the conference artifact

repository. The repository and package are licensed under Apache-2.0.


\begin{itemize}
\item \texttt{A1}: TempGNN reference model, TDP/DDTC/OATS tests, and real edge-stream
profiling.
\item \texttt{A2}: TempGNN plus independent paper-based MATG, ViTeGNN, and RTGA U280
forward paths, validated XRT hosts, xclbins, reports, and measurement workflow.
\item \texttt{A3}: Source-labeled paper-reference data and deterministic CSV/SVG
generator.
\end{itemize}

\begin{center}
\footnotesize
\begin{tabular}{@{}p{0.25\linewidth}p{0.25\linewidth}p{0.36\linewidth}@{}}
\toprule
Artifact & Contributions supported & Related paper elements \\
\midrule
\texttt{A1} & \texttt{C1}, \texttt{C2}, \texttt{C3} & TDP correctness and mechanism behavior; Fig.4(a), Fig.13, Fig.14 context \\
\addlinespace[2pt]
\texttt{A2} & \texttt{C2}, \texttt{C3}, bounded evidence for \texttt{C4} & U280 functional evidence and fresh Fig.11/Fig.12-shaped diagnostic tables \\
\addlinespace[2pt]
\texttt{A3} & Audit record for \texttt{C4} & Fig.2, Fig.4(a), Fig.9(b), and Fig.10-Fig.14 \\
\addlinespace[2pt]
\bottomrule
\end{tabular}
\normalsize
\end{center}

\subsection{Measurement Boundary}

The package includes current 21-CU TempGNN board logs and post-route reports,

plus historical U280 sanity/layout evidence. Artifact \texttt{A2} executes four

distinct xclbins and collects fresh U280 latency, total-board power, clock, checksum, and input

provenance rows. MATG, ViTeGNN, and RTGA are independently written,

paper-based forward-path reproductions, not the baseline authors' complete

source trees.


The fresh four-system path is deliberately marked as bounded and diagnostic.

It uses 8-dimensional Q10 kernels, deterministic stand-in weights, and

8,192-event real-data prefixes. The paper uses complete model configurations,

default 32-bit floating point, trained checkpoints, full evaluation streams,

and post-route Vivado power estimates. Artifact \texttt{A3} regenerates packaged

paper-reference figures from documented numeric inputs; it does not execute

the external baseline stacks. The repository therefore sets

\texttt{results\_reproduced\_eligible: false} and does not currently assert the Results

Reproduced badge.


\section{Part 2: Artifact Identification}

\section{Computational Artifact A1}

\subsection{Relation to Contributions}

\texttt{A1} is the inspectable software specification for TDP construction, temporal

dependency traversal, DDTC scheduling, PHLE packet reuse, and OATS overlap

handling. It supports \texttt{C1-C3} independently of FPGA availability and exposes

the counters used to reason about branch parallelism, packet overlap, reuse,

collisions, synchronization stalls, and off-chip traffic.


\subsection{Expected Results}

The unit suite must pass all 34 tests. The TDP implementation must agree with a

chronological reference update, fixture generation must be deterministic, and

artifact paths/provenance must pass their consistency checks. Optional Q14

profiling produces nonnegative, well-formed per-dataset/model statistics from

real temporal edge streams. These checks substantiate the functional behavior

of the execution mechanisms in \texttt{C1-C3}; they are not performance measurements

of the full paper system.


\subsection{Expected Reproduction Time (in Minutes)}

\begin{center}
\footnotesize
\begin{tabular}{@{}p{0.25\linewidth}p{0.25\linewidth}p{0.36\linewidth}@{}}
\toprule
Step & Time & Notes \\
\midrule
Setup & 5-10 & Create a Python environment and install \texttt{requirements.txt} \\
\addlinespace[2pt]
Execution & less than 2 & Unit tests; optional full Q14 download/profile takes 10-40 \\
\addlinespace[2pt]
Analysis & less than 2 & Inspect test output and generated CSV/Markdown summaries \\
\addlinespace[2pt]
\bottomrule
\end{tabular}
\normalsize
\end{center}

\subsection{Artifact Setup}

\subsubsection{Hardware}

Any Linux x86\_64 machine is sufficient. No FPGA is required for \texttt{A1}.


\subsubsection{Software}

\begin{itemize}
\item Python 3.10 or newer: \path{https://www.python.org/}.
\item GNU Make and Bash: \path{https://www.gnu.org/software/make/}.
\item NumPy 1.23 or newer for the complete software workflow.
\item The exact recorded environment is in \texttt{ENVIRONMENT.md}.
\end{itemize}

\subsubsection{Datasets and Inputs}

Unit tests generate deterministic fixtures. Optional real-edge profiling uses

TGL edge streams from \path{https://github.com/amazon-science/tgl}. The package also

contains fixed 8,192-event prefixes for WK, MC, RT, LM, WT, and GT under

\path{external/u280_dataset_samples/}, each with its source URL, selection rule, and

SHA256 metadata.


\subsubsection{Installation and Deployment}

\begin{Verbatim}[fontsize=\small,breaklines=true,breakanywhere=true]
python3 -m venv .venv
source .venv/bin/activate
pip install -r requirements.txt
\end{Verbatim}

\subsection{Artifact Evaluation}

The default dependency chain is:


\begin{Verbatim}[fontsize=\small,breaklines=true,breakanywhere=true]
T1: create environment
  -> T2: run 34 unit and artifact-contract tests
  -> T3: optionally fetch real edge streams
  -> T4: profile TDP/DDTC/OATS behavior
  -> T5: generate CSV and Markdown summaries
\end{Verbatim}

Commands:


\begin{Verbatim}[fontsize=\small,breaklines=true,breakanywhere=true]
python3 -m unittest discover -s tests
make data
make q14
\end{Verbatim}

Q14 defaults to WIKI, MOOC, and REDDIT; models are JODIE, TGAT, TGN, and APAN;

the target batch size, fanout, and depth are recorded in the output CSVs.


\subsection{Artifact Analysis}

The test command must finish with \texttt{Ran 34 tests} and \texttt{OK}. Optional Q14 outputs

are written to:


\begin{Verbatim}[fontsize=\small,breaklines=true,breakanywhere=true]
results/q14_real_tgl_edges/q14_dataset_model_summary.csv
results/q14_real_tgl_edges/q14_batches.csv
results/q14_real_tgl_edges/q14_summary.md
\end{Verbatim}

The summary reports packet hit/reuse rates, collisions, synchronization stalls,

memory traffic, and a 225 MHz cycle-model latency. It must not be interpreted

as a fresh board timing result.


\section{Computational Artifact A2}

\subsection{Relation to Contributions}

\texttt{A2} makes the TempGNN forward path and the paper-based MATG, ViTeGNN, and RTGA

mechanisms executable on the same Alveo U280. It provides functional U280

evidence for \texttt{C2-C3} and a controlled diagnostic comparison related to \texttt{C4}.

Each implementation has a separate source top, synthesis result, xclbin, and

provenance record. Missing or byte-identical xclbins cause preflight failure.


\subsection{Expected Results}

Preflight must report four distinct xclbin SHA256 values. Every hardware row

must pass the software golden output and repeat-consistency checks, contain a

real-input URL and hash, and record the xclbin link request, Vivado-connected

kernel clock, and post-route WNS/TNS.

With 3 repetitions, the full matrix contains 288 raw rows:


\begin{Verbatim}[fontsize=\small,breaklines=true,breakanywhere=true]
4 implementations x 6 datasets x 4 models x 3 repetitions
\end{Verbatim}

The packaged run \path{20260717T024537Z} satisfies this functional contract:
all 288 rows pass golden, repeat-consistency, and timing validation. Its fresh
all-workload averages are \texttt{1.296553 ms} for TempGNN and
\texttt{9.966618 ms} for MATG, corresponding to \textbf{7.8889x}
measured average speedup.

The workflow derives fresh Fig.11/Fig.12-shaped CSV/SVG tables and a numerical

tolerance report. A tolerance failure is retained as evidence and is not

replaced with packaged reference values. Because the configuration is not

paper-equivalent, even a numerical pass is diagnostic only.


\subsection{Expected Reproduction Time (in Minutes)}

\begin{center}
\footnotesize
\begin{tabular}{@{}p{0.25\linewidth}p{0.25\linewidth}p{0.36\linewidth}@{}}
\toprule
Step & Time & Notes \\
\midrule
Setup & 10-20 & Receive reviewer account, source XRT, verify U280 \\
\addlinespace[2pt]
Execution & 30-90 & Direct run of four packaged xclbins, 3 repetitions \\
\addlinespace[2pt]
Analysis & less than 5 & Automatic aggregation, figures, tolerance and provenance \\
\addlinespace[2pt]
Optional rebuild & 240-480 per xclbin & Not required for review and outside the direct-run path \\
\addlinespace[2pt]
\bottomrule
\end{tabular}
\normalsize
\end{center}

\subsection{Artifact Setup}

\subsubsection{Hardware}

\begin{itemize}
\item AMD/Xilinx Alveo U280, device \texttt{xcu280-fsvh2892-2L-e}.
\item Platform \path{xilinx_u280_gen3x16_xdma_1_202211_1}.
\item The authors provide a time-bounded account on a matching U280 through the
private SC submission channel. No credentials are stored in the artifact.
\end{itemize}

\subsubsection{Software}

\begin{itemize}
\item Ubuntu 22.04 x86\_64, Linux 5.15 in the recorded run.
\item XRT 2.16.204: \path{https://github.com/Xilinx/XRT}.
\item Vitis/Vivado 2023.2 is required only to rebuild:
\path{https://www.xilinx.com/support/download.html}.
\item GCC 11.4, GNU Make 4.3, Bash, Python 3.10 or newer.
\end{itemize}

\subsubsection{Datasets and Inputs}

The direct workflow uses bundled, timestamp-sorted 8,192-event prefixes of

Wikipedia, MOOC, Reddit, LastFM, WikiTalk, and GDELT. Each prefix has source and

hash metadata. Fixtures are generated for JODIE, TGN, TGAT, and APAN with

batch size 1,000. Synthetic inputs are permitted only for C-sim and are

rejected by the core workflow. The repository and frozen archive retain the

exact generated fixture metadata and goldens under

\path{results/generated_u280_comparison_fixtures/} for direct hash inspection.


\subsubsection{Installation and Deployment}

\begin{Verbatim}[fontsize=\small,breaklines=true,breakanywhere=true]
source /opt/xilinx/xrt/setup.sh
source /tools/Xilinx/Vitis/2023.2/settings64.sh  # rebuild only
make u280-core-preflight
\end{Verbatim}

The reviewer-facing files are under \path{artifacts/u280/}; source and mechanism

mapping are under \path{hardware/} and \path{hardware/baselines/}.


\subsection{Artifact Evaluation}

The direct-run dependency chain is:


\begin{Verbatim}[fontsize=\small,breaklines=true,breakanywhere=true]
T1: preflight hashes, files, source revisions, and four xclbins
  -> T2: generate deterministic real-prefix fixtures and software goldens
  -> T3: run each xclbin with its packaged validated XRT host
  -> T4: sample gated U280 power and record per-repetition raw rows
  -> T5: validate completeness, checksums, energy, and timing-closed clocks
  -> T6: aggregate rows and derive measured Fig.11/Fig.12 comparison tables
  -> T7: compare with source-labeled paper-reference values
\end{Verbatim}

Commands:


\begin{Verbatim}[fontsize=\small,breaklines=true,breakanywhere=true]
make ae-core-u280 U280_CORE_DEVICE=0 U280_CORE_REPETITIONS=3
\end{Verbatim}

The equivalent default-device wrapper is \path{bash scripts/run_all.sh u280-core}.

The checked parameters are in \path{configs/u280_core_reproduction.json}: each

design's requested timing-closed clock and CU count, no frequency rescaling,

batch size 1,000, six datasets, four models, three repetitions, 8-dimensional Q10 tensors, and bounded real-data

prefixes.


\subsection{Artifact Analysis}

Fresh results are written to a new timestamped directory:


\begin{Verbatim}[fontsize=\small,breaklines=true,breakanywhere=true]
results/reviewer_u280_runs/<run-id>/provenance.json
results/reviewer_u280_runs/<run-id>/raw/*/measurements.csv
results/reviewer_u280_runs/<run-id>/baselines_u280/
results/reviewer_u280_runs/<run-id>/derived_comparison_figures/
results/reviewer_u280_runs/<run-id>/verification.json
results/reviewer_u280_runs/<run-id>/verification.md
\end{Verbatim}

The packaged final run is \path{20260717T024537Z}; it contains 288 measured
rows and is the run cited by the generated AE report. All functional checks
pass. The Fig.11/Fig.12 paper-tolerance diagnostic remains FAIL and is
preserved rather than replacing fresh rows with paper-reference values.


Latency is XRT launch-to-completion kernel time. Power is total U280 board power

sampled only during the gated repeated-kernel window. Energy in millijoules is

\texttt{latency\_ms * power\_w}. The orchestrator rejects incomplete matrices, failed

goldens, inconsistent energy, synthetic core inputs, and unstable or unverified

achieved clocks. Each implementation's timing-closed clock is recorded without

frequency rescaling. Post-route timing, utilization, HLS, route, build, and xclbin metadata

evidence is staged with SHA256 records.


\section{Computational Artifact A3}

\subsection{Relation to Contributions}

\texttt{A3} preserves the numerical inputs used to plot the paper evaluation and makes

the transformation to CSV/SVG inspectable and deterministic. It is an audit

artifact for \texttt{C4}, not a substitute for executing \texttt{A2} or external CPU/GPU

baseline stacks.


\subsection{Expected Results}

The command regenerates nine paper-reference figures and matching CSV files.

Every input row is labeled as either an exact author-workbook cell or an

axis-calibrated recovery from an author vector export. No value is backsolved

from a reported mean or silently replaced by a fresh U280 row.


Expected summary values include:


\begin{Verbatim}[fontsize=\small,breaklines=true,breakanywhere=true]
TempGNN vs TGLite-CPU: 132.80x
TempGNN-G vs TGLite-CPU: 10.85x
Cascade vs TGLite-CPU: 5.22x
TempGNN vs MATG, explicit plotted AVG: 7.7889x
Paper prose for TempGNN vs MATG: 7.6x
Energy, Cascade/TempGNN, explicit plotted AVG: 33.545x
w/o DDTC normalized time: 3.08x
w/o OATS normalized time: 1.77x
\end{Verbatim}

\subsection{Expected Reproduction Time (in Minutes)}

\begin{center}
\footnotesize
\begin{tabular}{@{}p{0.25\linewidth}p{0.25\linewidth}p{0.36\linewidth}@{}}
\toprule
Step & Time & Notes \\
\midrule
Setup & less than 5 & Same Python environment as \texttt{A1} \\
\addlinespace[2pt]
Execution & less than 1 & Deterministic CSV/SVG generation \\
\addlinespace[2pt]
Analysis & 2-5 & Inspect manifest, source labels, and expected summaries \\
\addlinespace[2pt]
\bottomrule
\end{tabular}
\normalsize
\end{center}

\subsection{Artifact Setup}

\subsubsection{Hardware}

Any Linux x86\_64 machine; no FPGA or GPU is required.


\subsubsection{Software}

Python 3.10 or newer. Core CSV/SVG generation uses the Python standard library.


\subsubsection{Datasets and Inputs}

The 668 source-labeled numeric records are structured constants in

\path{tempgenn/paper_reference_data.py}. \path{reference_inputs/README.md} records source hashes,

workbook cell or vector geometry locators, uncertainty, and the Fig.11

plot/prose discrepancy. Deterministic CSV/SVG outputs are included under
\path{results/paper_reproduction/} for direct inspection and regenerated from
the code-embedded records.


\subsubsection{Installation and Deployment}

Use the environment prepared for \texttt{A1}; no additional deployment is required.


\subsection{Artifact Evaluation}

\begin{Verbatim}[fontsize=\small,breaklines=true,breakanywhere=true]
T1: read and validate the source-labeled code constants
  -> T2: group rows by figure, dataset, model, and solution
  -> T3: emit matching CSV and SVG files
  -> T4: write a manifest and combined source table
  -> T5: run consistency tests
\end{Verbatim}

Command:


\begin{Verbatim}[fontsize=\small,breaklines=true,breakanywhere=true]
python3 -m scripts.reproduce_paper_figures
\end{Verbatim}

\subsection{Artifact Analysis}

Outputs are under \path{results/paper_reproduction/}, including Fig.2, Fig.4(a),

Fig.9(b), and Fig.10-Fig.14 CSV/SVG pairs, \path{all_figure_data.csv}, and

\path{figure_data_manifest.csv}. The command also reconstructs the complete 668-row

\path{paper_figure_values.csv} with the migration SHA-256. The manifest identifies the source module and data

status for every figure. The explicit Fig.11 AVG bar and the rounded 7.6x prose

statement are both preserved rather than forced to agree.

\newpage
\appendixAE

Paper ID: \texttt{pap142}


This appendix evaluates the three computational artifacts defined in the AD:


\begin{itemize}
\item \texttt{A1}: TempGNN reference model, tests, and edge-stream profiling.
\item \texttt{A2}: four-system U280 forward-path execution and measurement workflow.
\item \texttt{A3}: source-labeled paper-reference CSV/SVG reconstruction.
\end{itemize}

The recommended reviewer order is \texttt{A1 -> A3 -> A2}. The first two paths do not

require FPGA access. Artifact \texttt{A2} uses an author-provided Alveo U280 account;

credentials are delivered only through the private SC submission channel.


\section{Evaluation Scope and Badge Boundary}

The artifact is prepared to support Artifacts Available and Artifacts

Evaluated-Functional. Apache-2.0 is included; a version-specific DOI must be

present by the August 25, 2026 artifact freeze.


The current package does not assert Results Reproduced. The fresh U280 path

uses bounded 8,192-event prefixes, 8-dimensional Q10 kernels, deterministic

stand-in weights, and \texttt{xbutil} total-board power. The paper uses complete model

configurations, default 32-bit floating point, trained checkpoints, full

evaluation streams, and post-route Vivado power estimates. The configuration

therefore records \texttt{results\_reproduced\_eligible: false}. Packaged reference rows

are never substituted for a failed fresh measurement.


Expected direct review time, excluding optional Vitis rebuilds:


\begin{center}
\footnotesize
\begin{tabular}{@{}p{0.34\linewidth}p{0.58\linewidth}@{}}
\toprule
Artifact & Setup, execution, and analysis time \\
\midrule
\texttt{A1} & 5-10 min setup; less than 2 min execution; less than 2 min analysis \\
\addlinespace[2pt]
\texttt{A2} & 10-20 min setup; 30-90 min execution; less than 5 min analysis \\
\addlinespace[2pt]
\texttt{A3} & less than 5 min setup; less than 1 min execution; 2-5 min analysis \\
\addlinespace[2pt]
\bottomrule
\end{tabular}
\normalsize
\end{center}

\section{Computational Artifact A1}

\subsection{Artifact Setup}

Use Linux x86\_64 with Python 3.10 or newer, GNU Make, and Bash:


\begin{Verbatim}[fontsize=\small,breaklines=true,breakanywhere=true]
git clone https://github.com/TempGNN-Program/TempGNN.git
cd TempGNN
python3 -m venv .venv
source .venv/bin/activate
pip install -r requirements.txt
\end{Verbatim}

The frozen package may be used instead of \texttt{git clone}:


\begin{Verbatim}[fontsize=\small,breaklines=true,breakanywhere=true]
tar -xzf pap142_tempgnn_sc26_ae_u280.tgz
cd pap142_tempgnn_sc26_ae
\end{Verbatim}

No FPGA is needed. Unit fixtures are generated locally. Optional Q14 profiling

downloads TGL edge streams; packaged U280-prefix samples are already under

\path{external/u280_dataset_samples/} with URL and SHA256 metadata.


\subsection{Artifact Execution}

The task dependency chain is:


\begin{Verbatim}[fontsize=\small,breaklines=true,breakanywhere=true]
A1-T1 environment
  -> A1-T2 unit tests for TDP/DDTC/OATS and artifact contracts
  -> A1-T3 optional real-edge acquisition
  -> A1-T4 optional edge-stream profiling
  -> A1-T5 CSV/Markdown summary generation
\end{Verbatim}

Run the required fast path:


\begin{Verbatim}[fontsize=\small,breaklines=true,breakanywhere=true]
python3 -m unittest discover -s tests
\end{Verbatim}

Expected terminal summary:


\begin{Verbatim}[fontsize=\small,breaklines=true,breakanywhere=true]
Ran 34 tests
OK
\end{Verbatim}

Run the optional real-edge path:


\begin{Verbatim}[fontsize=\small,breaklines=true,breakanywhere=true]
make data
make q14
\end{Verbatim}

The default Q14 parameters are WIKI/MOOC/REDDIT, four models, fanout 20, depth

2, and the batch parameters recorded in each output row. \texttt{make data} is skipped

when the corresponding edge files are already present.


\subsection{Artifact Analysis}

All 34 tests must pass. In particular, the recursive TDP state must agree with

chronological updates, fixtures must be deterministic, xclbin link requests

must agree with the Vivado-connected clock and nonnegative post-route WNS/TNS,

and duplicate xclbins must be rejected by preflight.


Optional Q14 outputs:


\begin{Verbatim}[fontsize=\small,breaklines=true,breakanywhere=true]
results/q14_real_tgl_edges/q14_dataset_model_summary.csv
results/q14_real_tgl_edges/q14_batches.csv
results/q14_real_tgl_edges/q14_summary.md
\end{Verbatim}

Inspect packet hit/reuse rates, collisions, stalls, and memory traffic. The

latency columns in this CPU path use a 225 MHz cycle model and are not fresh

U280 board timings. Passing \texttt{A1} demonstrates that the mechanisms supporting

contributions \texttt{C1-C3} are executable and internally checked.


\section{Computational Artifact A2}

\subsection{Artifact Setup}

Hardware and recorded software environment:


\begin{Verbatim}[fontsize=\small,breaklines=true,breakanywhere=true]
Board: AMD/Xilinx Alveo U280, xcu280-fsvh2892-2L-e
Platform: xilinx_u280_gen3x16_xdma_1_202211_1
OS: Ubuntu 22.04 x86_64
XRT: 2.16.204
Vitis/Vivado: 2023.2 (rebuild only)
GCC: 11.4.0
GNU Make: 4.3
Python: 3.10 or newer
\end{Verbatim}

The reviewer receives a dedicated, time-bounded account through the SC private

channel. After login:


\begin{Verbatim}[fontsize=\small,breaklines=true,breakanywhere=true]
source /opt/xilinx/xrt/setup.sh
cd /path/to/pap142_tempgnn_sc26_ae
xrt-smi examine
\end{Verbatim}

The four metadata-normalized artifacts are:


\begin{Verbatim}[fontsize=\small,breaklines=true,breakanywhere=true]
artifacts/u280/TempGNN/bin/tempgnn_forward_kernel.hw.xclbin
artifacts/u280/MATG/bin/matg_kernel.hw.xclbin
artifacts/u280/ViTeGNN/bin/vitegnn_kernel.hw.xclbin
artifacts/u280/RTGA/bin/rtga_kernel.hw.xclbin
\end{Verbatim}

Each validated XRT host and the baseline C-sim reference binaries are colocated

under the same \path{bin/} directories. TempGNN uses its 21-CU parallel host;

the baselines use the common single-CU host. Source, paper-mechanism mapping, and limitations

are in \path{hardware/baselines/README.md} and

\path{hardware/baselines/MECHANISM_MAP.md}.

The repository and frozen package also include the exact real-input fixture

metadata and goldens used by the packaged run under

\path{results/generated_u280_comparison_fixtures/}.


Optional rebuilds require:


\begin{Verbatim}[fontsize=\small,breaklines=true,breakanywhere=true]
source /tools/Xilinx/Vitis/2023.2/settings64.sh
\end{Verbatim}

Rebuilding is not required for evaluation and can take several hours per

xclbin. The direct-run path stays within the SC26 review budget.


\subsection{Artifact Execution}

The workflow is implemented by \path{scripts/run_u280_core_reproduction.py}:


\begin{Verbatim}[fontsize=\small,breaklines=true,breakanywhere=true]
A2-T1 preflight four source revisions, hosts, and distinct xclbin hashes
  -> A2-T2 generate fixtures from bundled real-data prefixes
  -> A2-T3 generate independent software goldens
  -> A2-T4 load and run each implementation on U280
  -> A2-T5 gate the repeated-kernel window and sample total board power
  -> A2-T6 validate checksums, repetitions, energy, provenance, and clocks
  -> A2-T7 aggregate raw rows and derive measured core comparison tables
  -> A2-T8 compare fresh tables with the paper-reference rows
\end{Verbatim}

Run preflight:


\begin{Verbatim}[fontsize=\small,breaklines=true,breakanywhere=true]
make u280-core-preflight
\end{Verbatim}

Expected: U280 four-implementation preflight PASS, followed by four different

xclbin SHA256 values.


Run the complete direct measurement:


\begin{Verbatim}[fontsize=\small,breaklines=true,breakanywhere=true]
make ae-core-u280 U280_CORE_DEVICE=0 U280_CORE_REPETITIONS=3
\end{Verbatim}

The default-device wrapper is \path{bash scripts/run_all.sh u280-core}.

Checked parameters are loaded from \path{configs/u280_core_reproduction.json}:


\begin{Verbatim}[fontsize=\small,breaklines=true,breakanywhere=true]
datasets: WK, MC, RT, LM, WT, GT
models: JODIE, TGN, TGAT, APAN
implementations: TempGNN, MATG, ViTeGNN, RTGA
real prefix: 8,192 events per dataset
target batch: 1,000
repetitions: 3
tensor path: 8-dimensional Q10
requested clocks: TempGNN 168 MHz; MATG/ViTeGNN/RTGA 225 MHz
achieved-clock comparison tolerance: 0.5 MHz
\end{Verbatim}

Each dataset/model pair is calibrated to a repeated-kernel measurement window.

The host performs a warmup, waits at a file gate, and then executes the measured

iterations while the harness samples \texttt{xbutil examine --report electrical}.


Optional one-fixture TempGNN sanity command:


\begin{Verbatim}[fontsize=\small,breaklines=true,breakanywhere=true]
make u280-run U280_DEVICE=0
\end{Verbatim}

Optional rebuild commands:


\begin{Verbatim}[fontsize=\small,breaklines=true,breakanywhere=true]
make u280-build \
  U280_PLATFORM=/path/to/xilinx_u280_gen3x16_xdma_1_202211_1.xpfm
make u280-baseline-build
make u280-stage-artifacts
\end{Verbatim}

Baseline links should be run serially on a 64 GB host. The staging command

redacts login, host, and absolute paths from textual evidence and the xclbin

\path{BUILD_METADATA}/\path{SYSTEM_METADATA} sections; it verifies that the FPGA

\texttt{BITSTREAM} section is byte identical before and after normalization.


\subsection{Artifact Analysis}

The run creates a new directory and never overwrites paper-reference data:


\begin{Verbatim}[fontsize=\small,breaklines=true,breakanywhere=true]
results/reviewer_u280_runs/<run-id>/provenance.json
results/reviewer_u280_runs/<run-id>/raw/TempGNN/measurements.csv
results/reviewer_u280_runs/<run-id>/raw/MATG/measurements.csv
results/reviewer_u280_runs/<run-id>/raw/ViTeGNN/measurements.csv
results/reviewer_u280_runs/<run-id>/raw/RTGA/measurements.csv
results/reviewer_u280_runs/<run-id>/baselines_u280/
results/reviewer_u280_runs/<run-id>/derived_comparison_figures/
results/reviewer_u280_runs/<run-id>/verification.json
results/reviewer_u280_runs/<run-id>/verification.md
\end{Verbatim}

The packaged final run is

\path{results/reviewer_u280_runs/20260717T024537Z/}; its four raw CSV files
contain 72 rows each (288 total), with zero golden, repeat, or timing failures.
The fresh all-workload averages are \texttt{1.296553 ms} for TempGNN
and \texttt{9.966618 ms} for MATG, yielding \textbf{7.8889x} measured
average speedup. The paper-tolerance diagnostic remains FAIL and is preserved.


For 3 repetitions, each implementation must contain 72 rows and the combined

matrix must contain 288 rows. Every row must report:


\begin{itemize}
\item \path{golden_validation=PASS} and \path{repeat_consistency=PASS};
\item a nonzero kernel and embedding checksum;
\item a real dataset source URL and 64-character input/golden hashes;
\item positive latency and power;
\item energy consistent with \texttt{latency\_ms * power\_w};
\item requested and xclbin-link-requested clocks, the Vivado-connected post-route
kernel clock, WNS/TNS, and \path{timing_met=PASS}.
\end{itemize}

The workflow refuses normalized comparison generation if an implementation's

post-route clock changes between rows or lacks timing closure. Each actual

clock remains in the raw rows and no frequency rescaling is applied.

\texttt{verification.md} reports the maximum relative error against

paper-reference Fig.11 and the U280 subset of Fig.12. The reviewer Make target

default one-command path preserves a tolerance FAIL while still completing a

functionally valid hardware run. \texttt{make ae-core-u280-strict} makes that

mismatch fail the command. It must not be converted to PASS by adjusting measured rows.


Passing \texttt{A2} establishes that four distinct implementations execute on the

same U280, produce stable software-checked outputs, and emit traceable raw

measurement data. The bounded reproduction scope prevents this path from independently

supporting a Results Reproduced claim for \texttt{C4}.


\section{Computational Artifact A3}

\subsection{Artifact Setup}

Use the Python environment created for \texttt{A1}. No FPGA, GPU, external dataset, or

network access is required. Inputs are:


\begin{Verbatim}[fontsize=\small,breaklines=true,breakanywhere=true]
tempgenn/paper_reference_data.py
reference_inputs/README.md
\end{Verbatim}

The source module has 668 structured records. Every record identifies its source class, locator, and

uncertainty. Original author workbooks and vector files are not distributed

because they contain author metadata; only the anonymized numeric extraction

and source hashes are embedded in code. The AE archive also includes
deterministic CSV/SVG outputs under \path{results/paper_reproduction/} for
direct inspection; the code-embedded records remain their authoritative source.


\subsection{Artifact Execution}

The dependency chain is:


\begin{Verbatim}[fontsize=\small,breaklines=true,breakanywhere=true]
A3-T1 validate the code-embedded records and source labels
  -> A3-T2 split records by figure
  -> A3-T3 generate matching CSV and SVG files
  -> A3-T4 generate the combined table and manifest
  -> A3-T5 run consistency tests
  -> A3-T6 regenerate the reviewer report
\end{Verbatim}

Commands:


\begin{Verbatim}[fontsize=\small,breaklines=true,breakanywhere=true]
python3 -m scripts.reproduce_paper_figures
python3 -m scripts.make_ae_report \
  --q14-summary results/q14_real_tgl_edges/q14_dataset_model_summary.csv \
  --board-json results/board_u280/summary.json \
  --out results/ae_report
\end{Verbatim}

Equivalent Make target:


\begin{Verbatim}[fontsize=\small,breaklines=true,breakanywhere=true]
make smoke
make report
\end{Verbatim}

\texttt{make smoke} includes the 34 unit and artifact-contract tests from \texttt{A1-A3}.


\subsection{Artifact Analysis}

The following outputs are included for direct inspection and regenerated in
place by the commands above:


\begin{Verbatim}[fontsize=\small,breaklines=true,breakanywhere=true]
results/paper_reproduction/paper_figure_values.csv
results/paper_reproduction/all_figure_data.csv
results/paper_reproduction/figure_data_manifest.csv
results/paper_reproduction/fig2_execution_breakdown.csv/.svg
results/paper_reproduction/fig4a_branch_parallelism_ratio.csv/.svg
results/paper_reproduction/fig9b_gpu_overhead_breakdown.csv/.svg
results/paper_reproduction/fig10_speedup_tglite_cpu.csv/.svg
results/paper_reproduction/fig11_speedup_matg.csv/.svg
results/paper_reproduction/fig12_energy_tempgnn.csv/.svg
results/paper_reproduction/fig13_ablation_time.csv/.svg
results/paper_reproduction/fig14a_batch_sensitivity.csv/.svg
results/paper_reproduction/fig14b_tdp_entries.csv/.svg
\end{Verbatim}

Expected summary checks:


\begin{Verbatim}[fontsize=\small,breaklines=true,breakanywhere=true]
TempGNN vs TGLite-CPU: 132.80x
TempGNN-G vs TGLite-CPU: 10.85x
Cascade vs TGLite-CPU: 5.22x
TempGNN vs MATG, explicit plotted AVG: 7.7889x
Paper prose for TempGNN vs MATG: 7.6x
Energy, Cascade/TempGNN, explicit plotted AVG: 33.545x
w/o DDTC normalized time: 3.08x
w/o OATS normalized time: 1.77x
\end{Verbatim}

The Fig.11 explicit AVG and the paper's rounded prose value are both preserved.

The generator must not alter individual values to force agreement. Generated

SVG titles contain \texttt{Paper-reference}, and the manifest identifies whether each

series came from exact workbook cells or axis-calibrated vector geometry.


Passing \texttt{A3} demonstrates deterministic reconstruction and provenance of the

code-embedded paper plotting records for \texttt{C4}. It does not demonstrate a fresh rerun

of MATG, ViTeGNN, RTGA, Cascade, TGLite-CPU, or the complete TempGNN paper

configuration.


\section{Badge Evidence Summary}

\begin{center}
\footnotesize
\begin{tabular}{@{}p{0.34\linewidth}p{0.58\linewidth}@{}}
\toprule
Badge & Current evidence \\
\midrule
Artifacts Available & Source, tests, inputs, CSV/SVG records, four metadata-normalized U280 xclbins, reports, and documentation; Apache-2.0 included and DOI required by freeze \\
\addlinespace[2pt]
Artifacts Evaluated-Functional & \texttt{A1} tests, baseline C-sim, four-xclbin preflight, board goldens, repeated outputs, power/latency rows, and post-route evidence \\
\addlinespace[2pt]
Results Reproduced & Not currently asserted because \texttt{A2} is a bounded mechanism-level comparison rather than the paper's full precision, checkpoints, streams, and power method \\
\addlinespace[2pt]
\bottomrule
\end{tabular}
\normalsize
\end{center}
\end{document}
